%%%%%%%%%%%%%%%%%%%%%%%%%%%%%%%%%%%%%%%% DECLARATION %%%%%%%%%%%%%%%%%%%%%%%%%%%%%%%%%%%%%%%%
\documentclass[12pt,czech]{article}
\usepackage{graphicx}
\usepackage[T1]{fontenc}
\usepackage{babel}
\usepackage[utf8]{inputenc} % set coding
\usepackage{czech}
\usepackage[legalpaper, portrait, margin=0.8in]{geometry}
\usepackage{fixltx2e}
\usepackage{hyperref}
\usepackage{textgreek}
\usepackage{textcomp}
\usepackage{adjustbox}
\usepackage{float}
\usepackage{tabto}
\usepackage{amsmath}
\newcommand\mytab{\tab \hspace{10mm}}
\usepackage{xcolor}

\author{
  Kulíšek, Vojtěch\\
  \texttt{xkulis03@stud.fit.vutbr.cz}
  \and
  Plevač, Lukáš\\
  \texttt{xpleva07@stud.fit.vutbr.cz}
  \and
  Šesták, Pavel\\
  \texttt{xsesta07@stud.fit.vutbr.cz}
  \and
  Koláríková, Marie\\
  \texttt{xkolar77@stud.fit.vutbr.cz}
}

\title{Calculator - User manual}
\date{12.11.2020}  


%%%%%%%%%%%%%%%%%%%%%%%%%%%%%%%%%%%%%%%% TITLE PAGE %%%%%%%%%%%%%%%%%%%%%%%%%%%%%%%%%%%%%%%%
\begin{document}
  \begin{titlepage}
		\begin{center}
			%\textsc{\LARGE Vysoké Učení Technické v Brně} \\[0.5cm]
			%{\LARGE Fakulta informačních technologií}

			\begin{figure}
				\center\includegraphics[scale=0.3]{images/logo.png}
			\end{figure}

			\vspace{3cm}

			\textsc{\LARGE 	Praktické aspekty vývoje software	} \\[0.5cm]
			\textsc{\LARGE 2020/2021} \\[3.5cm]

			\textbf{\LARGE Uživatelská příručka} \\
			\vspace{3cm}
		    \vspace{3cm}
		    
		    \Large{Kulíšek Vojtěch(xkulis03),\\
		    Koláríková Marie(xkolar77),\\
		    Plevač Lukáš(xpleva07),\\
		    Šesták Pavel(xsesta07)} \\[0.5cm]
		    
			
			\vspace{3cm}
                
		    \textsc{}{\LARGE Brno, \today}

		\end{center}
	\end{titlepage}

%%%%%%%%%%%%%%%%%%%%%%%%%%%%%%%%%%%%%%%% CONTENT TABLE %%%%%%%%%%%%%%%%%%%%%%%%%%%%%%%%%%%%%%%%
\begin{center}
	\tableofcontents %content generate
	\newpage %end of page
\end{center}
  
%%%%%%%%%%%%%%%%%%%%%%%%%%%%%%%%%%%%%%%% SECTIONS       %%%%%%%%%%%%%%%%%%%%%%%%%%%%%%%%%%%%%%%%

\section{Úvod}
Kalkulačka slouží k~počítání základních matematických operací. 
Dostupné matematické operace jsou sčítání, odčítání, násobení, dělení, mocnina, odmocnina, faktoriál a~modulo. Kalkulačka umožňuje zadávání komplexních výrazů.

\subsection{Precedence operátorů}
\begin{enumerate}
  \item faktoriál
  \item mocnina, odmocnina
  \item násobení, dělení, modulo
  \item sčítání, odčítání
\end{enumerate}


\newpage

\section{Instalace}
\smallskip
\subsection{Windows}
\subsubsection*{Zapnutí instalačního programu}
Přejděte do složky obsahující instalační program kalkulačky. \newline
Zapněte program \textbf{but\_fit\_ivs Setup~0.1.0.exe}.

\subsubsection*{Průběh instalace}
Po spuštění instalačního programu, počkejte, než se kalkulačka nainstaluje.
\begin{figure}[h]
\centering
\scalebox{0.5}{\includegraphics{images/instalace.png}}
\end{figure}

\subsubsection*{Dokončení instalace}
Po dokončení instalace se kalkulačka spustí.
\begin{figure}[h]
\centering
\scalebox{0.5}{\includegraphics{images/kalkulacka.png}}
\end{figure}

\subsection{Odinstalace}
Pokud chcete kalkulačku odinstalovat, zapněte \textbf{nastavení}, poté \textbf{aplikace}, v~seznamu najděte položku \textbf{but\_fit\_ivs~0.1.0} a~vyberte \textbf{odinstalovat}.
\begin{figure}[h]
\centering
\scalebox{0.5}{\includegraphics{images/odinstalace.png}}
\end{figure}

\newpage
\subsection{Linux}
V případě distribuce založené na Debianu pro instalaci \newline
použijeme příkaz: \textbf{sudo dpkg -i but\_fit\_ivs\_0.1.0\_amd64.deb}. \newline
Pro odinstalaci \textbf{sudo apt purge but-fit-ivs}.\newline
V případě jiné distribuce lze využit: \textbf{but\_fit\_ivs-0.1.0.AppImage}.


\subsection{Mac OS}
Pro instalaci na MacOS slouží soubor \textbf{but\_fit\_ivs-0.1.0-mac.zip}, který je na Mac zařízeních interpretován jako instalační soubor.

\newpage
\section{Pro vývojáře}
\subsection{Závislosti}
\begin{itemize}
    \item{wine Mono}
    \item{npm}
\end{itemize}
Po instalaci všech závislostí přejdeme do složky src. Příkazem \textbf{npm install} nainstalujeme moduly Node a závislosti pro projekt definované v \textbf{package.json}. Veškeré operace spojené s projektem jsou definované v rámci \textbf{src/Makefile} a je doporučené si tento makefile prostudovat před případnou prácí s projektem.

\subsection{Profiling}
Při profilingu jsme objevili problém s načítáním ze standardního vstupu s chybovou hláškou EACCESS permision denied i přes správně nastavená práva. Tato chyba se projevuje na určitých verzí Node. Pro profiling je nutný jinak nastavený package.json než pro spuštění a build. Tyto závislosti jsou nastaveny v makefile.

\newpage
\section{Práce s~kalkulačkou}
K~zadání čísel a~matematických operací je možné použít tlačítka aplikace nebo klávesnici. Před použitím klávesnice je nutné myší vybrat vstup dat, který se nachází v horní části aplikace. 

\subsection{Ovládání pomocí klávesnice}
Pomocí vkládání lze standardně vkládat čísla, závorky a základní operátory.
Jako separátor desetinného čísla lze využít jak tečku tak čárku. Speciální operátory jako dělení a násobení se automaticky upravují při vkládání. Pro znak násobení použijeme zástupný znak "*". Obdobně pro dělení použijeme znak "/", který se přeloží na znak "÷". Odmocninu lze zapsat pomocí příkazu "\textbackslash sqrt".

\begin{figure}[h]
\centering
\scalebox{0.5}{\includegraphics{images/zadavani.png}}
\caption{Příklad vstupu kalkulačky}
\end{figure}

 
\begin{figure}[h]
\centering
\scalebox{0.5}{\includegraphics{images/vysledek.png}}
\caption{Pro zobrazení výsledku slouží tlačítko, klávesa \textbf{"="}~nebo klávesa "\textbf{enter}".}
\end{figure}

\newpage


\begin{figure}[h]
\centering
\scalebox{0.5}{\includegraphics{images/AnsTlacitko.png}}
\caption{Tlačítko \textbf{Ans} umožní použít poslední výsledek k dalšímu výpočtu.}
\end{figure}


\begin{figure}[h]
\centering
\scalebox{0.5}{\includegraphics{images/smazat.png}}
\caption{Pokud potřebujete zadanou úlohu změnit, použijte tlačítko \textbf{C}, které vše smaže, nebo klávesu \textbf{delete}.}
\end{figure}

\section{Formát zadávání operací}
Klasické operace sčítání, odčítání, násobení, dělení se zadávají \newline
standardně ve formátu \textbf{operand1, operátor, operand2}.\\
\textbf{Například:} 7 - 4 = 3.

\subsection{Mocnina}
Mocnina se zadává ve formátu základ mocniny, symbol mocniny, exponent.\\
\textbf{Například:} 2**10 = 1024.

\subsection{Odmocnina}
Mocnina se zadává ve formátu \textbf{stupeň odmocniny, symbol odmocniny, číslo k odmocnění}. V případě stupně odmocniny rovné dvěma lze tento údaj vynechat. \\
\textbf{Například:} 2\sqrt16 = 4 = \sqrt16.

\subsection{Faktoriál}
Faktoriál se zadává ve formátu \textbf{číslo, symbol faktoriálu}.
\textbf{Například:} 5! = 120.
\end{document}
